% LaTeX Template For MATH 490 @ VCU
\documentclass{article}

\usepackage{hyperref}
\usepackage{amsmath}
\usepackage{amsthm}
\usepackage{amssymb}
\usepackage{geometry}
\geometry{a4paper, portrait, margin=2cm}


\begin{document}

\title{Project Report}

\author{Tom Barnes}

\maketitle              % typeset the title of the contribution

\section{Project Summary}
In this project, I formalised a proof of the fundamental theorem of arithmetic, which says that 
for any natural number $n > 1$, there exists prime numbers $p_1, p_2, \dots, p_k$ such that 
$n = p_1p_2\cdots p_k$, and that this set of primes is unique up to reordering, that is if 
$n = q_1q_2\cdots q_l$ for $q_j$ all prime, then $k = l$ and the $q_j$ are a rearrangement of 
the $p_i$.

\subsection{Proposal Extension}
In my original proposal, I stated that I would only be aiming to prove the existence half of 
the theorem, however, in my first implementation I proved that result and all the necessary 
important lemmas in under $200$ lines. I then refactored the code where possible and cut down the 
proof to around $110$ lines, and I had enough time so decided to prove the uniqueness half as well.

\subsection{Definitions}
I defined a prime factorisation as a structure 
$\verb|PrimeFactorisation (n : | \mathbb{N} \verb|) (L : List | \mathbb{N} \verb|)|$
with two fields
\begin{itemize}
  \item \verb|all_primes| : $\forall p \in$ \verb|L|, $p$ is prime
  \item \verb|prod_eq| : \verb|L.prod = n|
\end{itemize}
Apart from this, I used Mathlib's definition of primality: \verb|Nat.Prime n| as stated in the
project proposal, as well as \verb|Nat.Coprime m n|. 

\subsection{Results Used}
In the existence half of the proof, I used no major results from Mathlib
aside from trivial lemmas relating to arithmetic and list products. For the uniqueness half, for 
the sake of time, I used a few relevant results from Mathlib, the most important being:
\begin{itemize}
\item \verb|Nat.gcdA a b| and \verb|Nat.gcdB a b|
\newline One of my intermediate lemmas requires Bezout's theorem, which states: 
\[
\forall a, b \in \mathbb{N}, \exists\; x, y \in \mathbb{Z} \text{, s.t. } ax + by = \gcd(a, b) 
\]
These results from Mathlib respectively provide the $x$ and $y$ in this identity, and it would
have been a significant amount of work to prove this result from scratch.
\end{itemize}
I also used a few other results relating to primality and coprimality that would have been 
fairly trivial to prove:
\begin{itemize}
\item \verb|Nat.not_prime_iff_exists_dvd_ne| : $n$ is not prime $\iff \exists m, \; m \mid n, m \ne 1$ and $m \ne n$
\item \verb|Nat.Prime.coprime_iff_not_dvd| : if $p, n \in \mathbb{N}$ and $p$ is prime, 
$\gcd(p, n) = 1 \iff p \nmid n$ 
\item \verb|Nat.Prime.dvd_iff_eq| : if $p, n \in \mathbb{N}$, $n > 1$ and $p$ is prime, 
$n \mid p \iff n = p$
\item \verb|Nat.Prime.not_dvd_one| : if $p$ is prime, $p \nmid 1$
\item \verb|Nat.Prime.ne_one|, \verb|Nat.Prime.ne_zero| : if $p$ is prime, $p \ne 1$, $p \ne 0$

\end{itemize}

\newpage
\subsection{Results Proven}
Here I state the results I proved and a high-level outline of the proof where it 
is not trivial.
\subsubsection{Existence}

\begin{enumerate}
\item \verb|prime_factorisation_singleton| : if $p$ is prime, $\left[p\right]$ is a factorisation for $p$ 

\item \verb|empty_factorisation_is_one| : if a factorisation of $n$ is the empty list then $n = 1$

\item \verb|perm_is_factorisation| : if $L_1$ is a factorisation of $n$ and $L_2$ is a permutation
of $L_1$, then $L_2$ is also a factorisation of $n$
\newline \textit{Proof}: Use \verb|List.Perm.mem_iff| and \verb|List.Perm.prod_eq|.

\item \verb|factorisation_cons| : if $L=\left[p_1,\dots,p_k\right]$ is a factorisation of $n$ and 
$p$ is prime, then $\left[p, p_1, \dots, p_k\right]$ is a factorisation of $n\cdot p$
\newline \textit{Proof}: Use \verb|List.prod_cons| and \verb|List.mem_cons|.

\item \verb|factorisation_append| : if $L_1=\left[p_1,\dots,p_k\right], L_2=\left[q_1,\dots,q_m\right]$ are factorisations of $n_1$ and $n_2$ respectively, then $\left[p_1,\dots,p_k,q_1,\dots,q_m\right]$ is a factorisation of $n_1\cdot n_2$
\newline \textit{Proof}: Use induction on the length of $L_2$. For the base case, use \verb|empty_factorisation_is_one|, and simplify. For the induction step, proceed as follows:
\begin{itemize}
  \item show the tail of $L_2$ is a factorisation for the product of it's elements 
  \item use the induction hypothesis to combine $L_1$ with the tail of $L_2$
  \item use \verb|factorisation_cons| to add back the head of $L_2$
  \item show this is a permutation of $L_1$ combined with $L_2$ using \verb|List.perm_append_comm_assoc|
  \item use \verb|perm_is_factorisation| to conclude
\end{itemize}

\item \verb|not_prime_exists_factors| : if $n \ge 2$ is composite then it has factors $a, b$ with $1 < a, b < n$ 
\newline \textit{Proof}: Apply \verb|Nat.not_prime_iff_exists_dvd_ne| to obtain a factor $m$ of 
$n$ that is not equal to $1$ or $n$, then use \verb|grind|, \verb|refine| and \verb|by_cases| 
tactics to show that $m$ and $\frac{n}{m}$ are both strictly between $1$ and $n$.

\item \verb|exists_factorisation| : every natural number larger than $1$ has a prime factorisation
\newline \textit{Proof}: Use strong induction on $n$. In the case $n$ is a prime, apply 
\verb|prime_factorisation_singleton|. In the other case, apply \verb|not_prime_exists_factors| to
obtain nontrivial factors $a$ and $b$. Apply the induction hypothesis to obtain factorisations
for these, then apply \verb|factorisation_append| to finish the proof.

\end{enumerate}

\subsubsection{Uniqueness}

\begin{enumerate}

\item \verb|coprime_dvd_prod| : if $n \mid ab$ and $\gcd(a, n) = 1$ then $n \mid b$
\newline \textit{Proof}: Apply Bezout's lemma, multiply by $b$ and show $n$ divides both sides
of the resulting expression using Mathlib's \verb|Int.dvd_*| results.

\item \verb|prime_dvd_prod| : if $p$ is prime and $p \mid ab$ then $p \mid a$ or $p \mid b$
\newline \textit{Proof}: Trivial if $p \mid a$, otherwise use \verb|Nat.Prime.coprime_iff_not_dvd| 
to show $gcd(a, p) = 1$ then apply \verb|coprime_dvd_prod|.

\item \verb|prime_dvd_prod_dvd_factor| : if $p$ is prime and $p \mid n$, then $p$ divides a prime 
factor of $n$
\newline \textit{Proof}: Take a factorisation $L$ of $n$ and use induction on it's length. 
Use \verb|empty_factorisation_is_one| and \verb|Nat.Prime.not_dvd_one| to show $L$ is nonempty. 
Then if $p$ divides the product of tail of $L$, show it is a factorisation of it's product and
apply the induction hypothesis. Otherwise $p \mid n$ so apply \verb|prime_dvd_prod| to show $p$
must divide the head of $L$.

\item \verb|prime_factor_mem_factorisation| : if $L_1, L_2$ are factorisations for $n$ then
$p\in L_1 \implies p\in L_2$
\newline \textit{Proof}: Apply \verb|prime_dvd_prod_dvd_factor| to obtain $q \in L_2$ with 
$p \mid q$, then \verb|Nat.Prime.dvd_iff_eq|.

\item \verb|factorisation_remove_prime| : if $L$ is a factorisation for $n$ and $p \in L$ then
removing one copy of $p$ from $L$ is a factorisation for $n/p$
\newline \textit{Proof}: Split $L$ into a head, $L_1$, singleton containing $p$ and tail $L_2$, 
and apply induction. If $p$ is the first element of $L_2$, take $L = L_2$, otherwise apply the 
induction hypothesis to split $L_2$ further. Apply \verb|factorisation_cons| and simplify.

\item \verb|unique_factorisation| : if $L_1, L_2$ are both factorisations for $n$ then $L_2$ must
be a permutation of $L_1$
\newline \textit{Proof}: Apply induction on $L_1$. In the case it's empty, show $n=1$ by 
\verb|empty_factorisation_is_one| and that $L_2$ must also be empty. Otherwise, take
$p \in L_1$ and apply \verb|prime_factor_mem_factorisation| to get $p \in L_2$. Define $L'_1$ and
$L'_2$ as $L_1$ and $L_2$ respectively with one copy of $p$ removed. Apply induction on the 
shorter lists to show they are permutations. Then add back $p$ and show $L_1$ and $L_2$ are also
permutations.

\end{enumerate}

\section{Insights}
Since this is an elementary number theory result I have seen before many times, I didn't learn 
anything new of note regarding this proof, however I learnt a lot about Lean, including:
\begin{itemize}
\item \textbf{Tactics such as} \verb|refine, grind, simpa, obtain|
\item \textbf{Proof by induction including strong induction and generalizing on variables}
\item \textbf{Adhering to Mathlib conventions and Lean coding style}
\end{itemize}

\section{A.I. Usage}
For this project I used AI in the following ways:
\begin{itemize}
\item \textbf{Project structure:}
I wrote down a proof of the theorem then used ChatGPT to advise me on which lemmas would be suitable to prove, given the time constraints and my limited knowledge of lean.

\item \textbf{Mathlib theorem finding and examples:} 
I frequently used Google AI overview to find the relevant theorem on natural numbers or lists by 
describing what I wanted to achieve. I also used the AI overview to give me examples of using 
tactics such as \verb|refine| and \verb|induction generalizing|.

\end{itemize}

\end{document}
